\documentclass[../../../../main.tex]{subfiles}
\begin{document}

\begin{flushright}
\footnotesize\textit{【自動文字起こし・要確認】}
\end{flushright}

\noindent
[\textbf{解}] 答案を書き直し、修正部分には \underline{\hspace{0.8cm}} を下に引く。

\vspace{0.3cm}

\noindent
$(1)\times y - (2)\times x \implies (y-\underline{x})z(z+1)=0$

\noindent
から (イ) $z=0$, (ロ) $z=-1$, (ハ) $y=x$

\vspace{0.3cm}

\noindent
(イ)の時、(1),(2),(3)から $(x,y,z)=(0,0,0)$

\noindent
(ロ)の時、(1),(2)から $x=y=0$ だが、これは \underline{(3)を満たさず不適}

\noindent
(ハ)の時、(3)から $z^2(z+y+2)=0$, (イ)で $z=0$ は考えたので $z \neq 0$ として $z+y+2=0 \iff z=-(y+2)$ だから (2)に代入
\[
-y(y+2) = y + (y+2)^2 - (y+2)
\]
\[
y^2 + 3y + 1 = 0
\]
\[
y = \frac{-3 \mp \sqrt{5}}{2}
\]
$z=-(y+2)$ に代入して
\[
z = - \underline{\frac{1 \mp \sqrt{5}}{2} = \frac{-1 \pm \sqrt{5}}{2}}
\]
だから, $x = \frac{-3 \mp \sqrt{5}}{2}$

\vspace{0.3cm}

\noindent
以上から
\[
(x,y,z) = (0,0,0), \ \underline{\left( \frac{-3 \mp \sqrt{5}}{2}, \frac{-3 \mp \sqrt{5}}{2}, \frac{-1 \pm \sqrt{5}}{2} \right)}
\]

\newpage

\end{document}
